\SourceLeaf{451}
\par\Needspace{8\baselineskip}\bigskip
\pdfbookmark[3]{Darstellung, 4}{pallas-II-article4}
\begin{center}4.\end{center}
\noindent\hspace*{\parindent}Es bleibt noch, den Zusammenhang zwischen unseren Konstanten und den Elementen herzuleiten, die die Astronomen gewöhnlich verwenden. Zu diesem Zweck beziehen wir die folgenden neun Richtungen auf eine Kugeloberfläche von beliebigem Radius, deren Mittelpunkt mit dem Mittelpunkt der Sonne zusammenfällt. Es sind die drei Koordinatenachsen, die von der Sonne zum Planeten gezogene Gerade, eine in der Bahnebene auf ihr senkrechte Gerade, die auf dieser Ebene senkrechte Gerade, die zum Perihel gerichtete Apsidenlinie, die in der Bahnebene auf ihr senkrechte Gerade und der Schnitt der Ebene der \(x,y\) mit der Bahnebene. Wir bezeichnen die diesen Richtungen entsprechenden Punkte der Kugeloberfläche mit \(X,Y,Z,L,\Lambda,M,P,\Pi,N\). Somit sind \(X,Y,Z\) die Pole der Großkreise, die die drei Grundebenen darstellen. Wir nehmen an, dass diese Pole auf der Seite gewählt sind, auf der die Koordinaten als positiv gelten. Es sei \(L\) der heliozentrische Ort des Planeten, \(\Lambda\) der Punkt des Großkreises der Bahn, der \(L\) um \(90^{0}\) voraus ist, \(P\) das Perihel und \(\Pi\) der Punkt, in den \(\Lambda\) übergeht, wenn \(L\) in \(P\) übergeht. Es sei \(N\) der aufsteigende Knoten der Bahn auf dem Großkreis, dessen Pol \(Z\) ist. Schließlich sei \(M\) einer der Bahnpole. Um jede Mehrdeutigkeit auszuschließen, nehmen wir an, dass die Punkte \(P,\Pi,M\) auf der Kugel in derselben Reihenfolge angeordnet sind wie die Punkte \(X,Y,Z\). Dazu muss man \(M\) auf derselben Seite des Großkreises \(XY\) wie \(Z\) oder auf der gegenüberliegenden Seite wählen, je nachdem, ob die auf den Großkreis \(XY\) projizierte heliozentrische Bewegung des Planeten von \(X\) nach \(Y\) oder in der entgegengesetzten Richtung verläuft. Im ersten Fall heißt die Bewegung \textit{rechtläufig}, im zweiten \textit{rückläufig}. Um diese beiden Fälle analytisch zu unterscheiden, sei \(\beta\) der Abstand von \(L\) zum Großkreis \(XY\), oder genauer \(\beta=90^{0}-LZ\), und \(\lambda\) die Länge der Projektion von \(L\) auf diesen Großkreis, von \(X\) in der Richtung von \(X\) nach \(Y\) gezählt. Dann gilt
\[
\begin{aligned}
x&=r\cos\beta\cos\lambda\\
y&=r\cos\beta\sin\lambda\\
z&=r\sin\beta.
\end{aligned}
\]
\par Hieraus folgt \(x\mathrm{d}y-y\mathrm{d}x=rr\cos\beta^2\mathbin{.}\mathrm{d}\lambda\). Somit ist bei \SourceLeaf{452}rechtläufiger Bewegung \(C\) positiv, bei rückläufiger Bewegung negativ. Also haben \(C\) und \(\cos ZM\) stets dasselbe Vorzeichen. Nun gilt
\[
\begin{array}{rl}
27)\quad&x=r\cos XL\\[4pt]
28)\quad&y=r\cos YL\\[4pt]
29)\quad&z=r\cos ZL
\end{array}
\]
und nach einem bekannten Satz
\[
\begin{aligned}
&\cos XL\mathbin{.}\cos XM+\cos YL\mathbin{.}\cos YM+\cos ZL\mathbin{.}\cos ZM\\
&\qquad=\cos LM=0.
\end{aligned}
\]
\par Folglich gilt
\[
x\cos XM+y\cos YM+z\cos ZM=0
\]
und da diese Gleichung offenbar mit Gleichung 7 identisch sein muss, müssen die Größen \(A,B,C\) proportional zu \(\cos XM,\cos YM,\cos ZM\) sein. Nach einem anderen bekannten Satz gilt \(\cos XM^2+\cos YM^2+\cos ZM^2=1\). Wir haben soeben gezeigt, dass \(\cos ZM\) und \(C\) notwendig dasselbe Vorzeichen haben. Da außerdem \(AA+BB+CC=mp\) gilt, müssen wir schließen
\[
\begin{array}{rl}
30)\quad&A=\cos XM\mathbin{.}\BareRoot{mp}\\[4pt]
31)\quad&B=\cos YM\mathbin{.}\BareRoot{mp}\\[4pt]
32)\quad&C=\cos ZM\mathbin{.}\BareRoot{mp}.
\end{array}
\]
